\chapter{Justificación}

La propuesta se justifica por la necesidad de representar el estado del mercado, por la disponibilidad de datos verificables y por su viabilidad computacional.

\paragraph{Fundamento teórico.} Una intensidad con parámetros fijos no representa la dependencia temporal del flujo de órdenes. \citet{BacryHawkes2015} revisan evidencia de endogeneidad y agrupamiento en aplicaciones financieras de procesos de Hawkes. \citet{ContKukanovStoikov2014} documentan la relación entre desequilibrio del flujo y cambios de precio de corto plazo en acciones. \citet{MorariuPatrichiPakkanen2022} muestran que la excitación del flujo depende de variables de estado, como spread y desequilibrio. Estos resultados motivan una intensidad que combine memoria, periodicidad y covariables observables, sin presuponer que una mejor predicción produzca rentabilidad.

\paragraph{Pertinencia empírica.} BTC/USDT en Binance Spot permite estudiar un mercado continuo con datos públicos de alta frecuencia. \citet{BarucciEtAl2022} documentan patrones intradía y la importancia de los mercados denominados en USDT en la formación de precios de criptoactivos; \citet{FabreToke2025} aplican estimación Hawkes no paramétrica a su microestructura. La captura propia permite reconstruir el libro, verificar secuencias y conservar las transacciones originales. La evaluación compara así modelos sobre una base común y auditable.

\paragraph{Viabilidad.} Los métodos de árboles son adecuados para comparar relaciones no lineales en datos tabulares de dimensión moderada \citep{GrinsztajnEtAl2022}. El trabajo contrasta gradient boosting con random forest y elige el estimador según desempeño temporal y costo de cómputo. La combinación de esa estimación con una curva de distancia interpretable permite integrar el pronóstico en el control de inventario y evaluar por separado su calidad estadística y su utilidad económica.
